% GENERATED FILE. Edit canonical lessons, then run npm run generate:books.
% canonical-source-hash: fnv1a64:80a047b2a5ca3fe0
% canonical-lessons: LA-C323-pario, LA-C323-partus, LA-C323-obstetrix, LA-C323-nutrio, LA-C323-nutrix

\chapter{Īnfāns nātus est}
\label{ch:la-infans-natus-est}
\hlchaptermodality{323}

\begin{chapteropening}
\textbf{By the end of this chapter:} \emph{I can talk about a birth: to give birth, the birth, the midwife, to suckle and the nurse.}

Complete the last lesson of chapter 323: I can talk about a birth: to give birth, the birth, the midwife, to suckle and the nurse.
\end{chapteropening}

\section[pariō, parere]{pariō, parere — to give birth to; to produce}

\label{lesson:LA-C323-pario}



\begin{quote}
Before the new one: say the Latin for to perish, then the Latin for a guard.
\end{quote}

\subsection*{You'll want to know: pariō, parere}
\textbf{pariō, parere} — \textquotedblleft{}to give birth to; to produce\textquotedblright{}.

\begin{quote}
Uxor eius fīliam peperit.
\end{quote}

\begin{quote}
\emph{His wife gave birth to a daughter.}
\end{quote}

\emph{pariō, parere, peperī, partum}, short \emph{a}. Not \emph{pāreō, pārēre}, 'to obey' (taught). English \emph{parent} comes from its family.

\begin{cousinweb}
\textbf{pariō} is inherited from a Proto-Indo-European root meaning \textquotedblleft{}to bring forth\textquotedblright{} (EtymDB). English \emph{parent} comes from its participle \emph{parēns}.
\end{cousinweb}

\begin{grammarlens}[title={What you've built}]
One more word for telling, explaining and describing at B1.
\end{grammarlens}

\subsection*{Your turn}
\begin{itemize}
\raggedright
  \item \emph{Say it:} \emph{pariō, parere}
  \item \emph{Say it:} \emph{pariō, parere}, once more
  \item \emph{Say it:} a sentence of your own with \emph{pariō, parere}
  \item \emph{Recall:} say the Latin for to perish, then the Latin for a guard, then say \emph{pariō, parere} again
\end{itemize}

\begin{tcolorbox}[breakable,colback=teal!4,colframe=teal!35!black,title={Before you move on}]
What does pariō, parere mean? (\textquotedblleft{}To give birth to; to produce\textquotedblright{}.) Say it once more.
\end{tcolorbox}
\section[partus, partūs]{partus, partūs — a birth; offspring}

\label{lesson:LA-C323-partus}



\begin{quote}
Before the new one: say the Latin for a guard, then the Latin for to give birth to.
\end{quote}

\subsection*{You'll want to know: partus, partūs}
\textbf{partus, partūs} — \textquotedblleft{}a birth; offspring\textquotedblright{}.

\begin{quote}
Partus facilis erat, et māter laeta.
\end{quote}

\begin{quote}
\emph{The birth was easy, and the mother happy.}
\end{quote}

Masculine, 4th declension, from \emph{pariō}.

\begin{cousinweb}
\textbf{partus} is built on \emph{pariō}, \textquotedblleft{}to give birth\textquotedblright{} (EtymDB): a bearing, and what is born.
\end{cousinweb}

\begin{grammarlens}[title={What you've built}]
One more word for telling, explaining and describing at B1.
\end{grammarlens}

\subsection*{Your turn}
\begin{itemize}
\raggedright
  \item \emph{Say it:} \emph{partus, partūs}
  \item \emph{Say it:} \emph{partus, partūs}, once more
  \item \emph{Say it:} a sentence of your own with \emph{partus, partūs}
  \item \emph{Recall:} say the Latin for a guard, then the Latin for to give birth to, then say \emph{partus, partūs} again
\end{itemize}

\begin{tcolorbox}[breakable,colback=teal!4,colframe=teal!35!black,title={Before you move on}]
What does partus, partūs mean? (\textquotedblleft{}A birth; offspring\textquotedblright{}.) Say it once more.
\end{tcolorbox}
\section[obstetrīx, obstetrīcis]{obstetrīx, obstetrīcis — a midwife}

\label{lesson:LA-C323-obstetrix}



\begin{quote}
Before the new one: say the Latin for to give birth to, then the Latin for a birth.
\end{quote}

\subsection*{You'll want to know: obstetrīx, obstetrīcis}
\textbf{obstetrīx, obstetrīcis} — \textquotedblleft{}a midwife\textquotedblright{}.

\begin{quote}
Obstetrīx māne ad domum vēnit.
\end{quote}

\begin{quote}
\emph{The midwife came to the house in the morning.}
\end{quote}

Feminine, 3rd declension: she stands by the mother. English \emph{obstetrics}.

\begin{cousinweb}
\textbf{obstetrīx} is built on \emph{obstō}, \textquotedblleft{}to stand before\textquotedblright{} (EtymDB), with the feminine ending \emph{-trīx}: she who stands by. English \emph{obstetrics} comes from it.
\end{cousinweb}

\begin{grammarlens}[title={What you've built}]
One more word for telling, explaining and describing at B1.
\end{grammarlens}

\subsection*{Your turn}
\begin{itemize}
\raggedright
  \item \emph{Say it:} \emph{obstetrīx, obstetrīcis}
  \item \emph{Say it:} \emph{obstetrīx, obstetrīcis}, once more
  \item \emph{Say it:} a sentence of your own with \emph{obstetrīx, obstetrīcis}
  \item \emph{Recall:} say the Latin for to give birth to, then the Latin for a birth, then say \emph{obstetrīx, obstetrīcis} again
\end{itemize}

\begin{tcolorbox}[breakable,colback=teal!4,colframe=teal!35!black,title={Before you move on}]
What does obstetrīx, obstetrīcis mean? (\textquotedblleft{}A midwife\textquotedblright{}.) Say it once more.
\end{tcolorbox}
\section[nūtriō, nūtrīre]{nūtriō, nūtrīre — to suckle, to feed, to rear}

\label{lesson:LA-C323-nutrio}



\begin{quote}
Before the new one: say the Latin for a birth, then the Latin for a midwife.
\end{quote}

\subsection*{You'll want to know: nūtriō, nūtrīre}
\textbf{nūtriō, nūtrīre} — \textquotedblleft{}to suckle, to feed, to rear\textquotedblright{}.

\begin{quote}
Māter ipsa īnfantem nūtrīvit.
\end{quote}

\begin{quote}
\emph{The mother fed the baby herself.}
\end{quote}

\emph{nūtriō, nūtrīre, nūtrīvī, nūtrītum}. English \emph{nourish} and \emph{nutrition} come from it.

\begin{cousinweb}
\textbf{nūtriō} is an inherited verb meaning first \textquotedblleft{}to suckle\textquotedblright{} (EtymDB). English \emph{nourish} comes from it through French, and \emph{nutrition} directly.
\end{cousinweb}

\begin{grammarlens}[title={What you've built}]
One more word for telling, explaining and describing at B1.
\end{grammarlens}

\subsection*{Your turn}
\begin{itemize}
\raggedright
  \item \emph{Say it:} \emph{nūtriō, nūtrīre}
  \item \emph{Say it:} \emph{nūtriō, nūtrīre}, once more
  \item \emph{Say it:} a sentence of your own with \emph{nūtriō, nūtrīre}
  \item \emph{Recall:} say the Latin for a birth, then the Latin for a midwife, then say \emph{nūtriō, nūtrīre} again
\end{itemize}

\begin{tcolorbox}[breakable,colback=teal!4,colframe=teal!35!black,title={Before you move on}]
What does nūtriō, nūtrīre mean? (\textquotedblleft{}To suckle, to feed, to rear\textquotedblright{}.) Say it once more.
\end{tcolorbox}
\section[nūtrīx, nūtrīcis]{nūtrīx, nūtrīcis — a nurse, a wet nurse}

\label{lesson:LA-C323-nutrix}



\begin{quote}
Before the new one: say the Latin for a midwife, then the Latin for to suckle.
\end{quote}

\subsection*{You'll want to know: nūtrīx, nūtrīcis}
\textbf{nūtrīx, nūtrīcis} — \textquotedblleft{}a nurse, a wet nurse\textquotedblright{}.

\begin{quote}
Nūtrīx īnfantem in bracchiīs tenēbat.
\end{quote}

\begin{quote}
\emph{The nurse was holding the baby in her arms.}
\end{quote}

Feminine, 3rd declension, from \emph{nūtriō}.

\begin{cousinweb}
\textbf{nūtrīx} is built on \emph{nūtriō}, \textquotedblleft{}to suckle, to feed\textquotedblright{} (EtymDB), with the feminine ending \emph{-trīx}.
\end{cousinweb}

\begin{grammarlens}[title={What you've built}]
One more word for telling, explaining and describing at B1.
\end{grammarlens}

\subsection*{Your turn}
\begin{itemize}
\raggedright
  \item \emph{Say it:} \emph{nūtrīx, nūtrīcis}
  \item \emph{Say it:} \emph{nūtrīx, nūtrīcis}, once more
  \item \emph{Say it:} a sentence of your own with \emph{nūtrīx, nūtrīcis}
  \item \emph{Recall:} say the Latin for a midwife, then the Latin for to suckle, then say \emph{nūtrīx, nūtrīcis} again
\end{itemize}

\begin{tcolorbox}[breakable,colback=teal!4,colframe=teal!35!black,title={Before you move on}]
What does nūtrīx, nūtrīcis mean? (\textquotedblleft{}A nurse, a wet nurse\textquotedblright{}.) Say it once more.
\end{tcolorbox}
